\clearpage
\begin{explanationbox}
\textbf{\textcolor{CExplanation}{一、先说人话：这个规律为什么有用？}}

一眼看到 $y=\left(\frac13\right)^{x^2-2x}$，同时出现平方和指数，
很难靠“看整个式子”判断增减。
但计算它时，实际上要做两件事：先算 $t=x^2-2x$，
再算 $y=\left(\frac13\right)^t$。
因此我们可以\textbf{把难题拆成两个简单问题}：
\begin{enumerate}[leftmargin=*, itemsep=0.25em]
\item $x$ 增大时，$t$ 增大还是减小？
\item $t$ 增大时，$y$ 增大还是减小？
\end{enumerate}
把这两个答案接起来，就知道 $x$ 增大时 $y$ 怎样变化。
这是复合函数要研究单调性的真正动机：\textbf{用已知的简单函数，判断复杂函数。}

\medskip
\textbf{\textcolor{CExplanation}{二、用数字亲眼看见“减、减、增”}}

取 $g(x)=10-x$，$f(t)=10-t$。让 $x$ 从 $2$ 变成 $5$：
\[
\begin{array}{c|c|c}
 x & t=10-x & y=10-t\\ \hline
 2 & 8 & 2\\
 5 & 5 & 5
\end{array}
\]
$x$ 增大了；中间的 $t$ 反而减小；外层看见更小的 $t$，
算出的 $y$ 反而更大。因此两层都递减，整体递增。
\textbf{注意：这组数字是帮助理解的例子，严格证明要看后面的任意两点比较。}

\medskip
\textbf{\textcolor{CExplanation}{三、比背口诀更可靠的一句话}}

\textbf{递增的操作保留原来的大小顺序；递减的操作交换原来的大小顺序。}
比如 $3<7$：
\[
 3+2<7+2\quad(\text{大小顺序不变}),
 \qquad -3>-7\quad(\text{大小顺序反过来}).
\]
复合函数只不过是连续做两次这样的比较。
从 $x_1<x_2$ 出发，经过两次处理：
\[
\begin{array}{c|c|c}
 \text{第一步} & \text{第二步} & \text{最终比较}\\ \hline
 \text{不反向} & \text{不反向} & y_1<y_2\\
 \text{反向} & \text{反向} & y_1<y_2\\
 \text{不反向} & \text{反向} & y_1>y_2\\
 \text{反向} & \text{不反向} & y_1>y_2
\end{array}
\]
这就把四种组合一次讲清楚了。

\medskip
\end{explanationbox}
\clearpage
\begin{explanationbox}
\textbf{\textcolor{CExplanation}{四、最关键的一步：先找外层实际接收什么数}}
\begin{center}\begin{tikzpicture}[font=\small,>=stealth,line join=round]
\draw[gray!40,rounded corners=3pt] (0,0) rectangle (14.2,9.4);
\node at (7.1,8.8) {例5：内层跨过零，外层平方要分成两边判断};

\begin{scope}[shift={(2.5,2.3)},x=1.2cm,y=.85cm]
\draw[->,gray] (-1.5,0)--(3.8,0) node[right]{$x$};\draw[->,gray] (0,-.3)--(0,5.3) node[above]{$H$};
\draw[blue,very thick,domain=-1.2:3.2,samples=140] plot(\x,{(\x-1)^2});\draw[gray,dotted] (1,0)--(1,4.8);\fill[orange!85!black] (1,0) circle(2.4pt);\node[below=6pt] at (1,0){$1$};
\end{scope}
\node[align=left,anchor=west] at (8.3,5){$t=x-1$：内层一直增\\[12pt]$t<0$ 时平方递减\\[9pt]$t>0$ 时平方递增\\[14pt]$x\le1$：整体减\\[9pt]$x\ge1$：整体增};

\node[text width=12.8cm,align=center] at (7.1,.65) {外层转折发生在 $t=0$，换回原来的自变量位置就是 $x=1$。};
\end{tikzpicture}\end{center}



以 $y=(x-1)^2$ 为例，写成 $g(x)=x-1$、$f(t)=t^2$。
内层 $g$ 在整个实数范围内递增，但外层 $f(t)=t^2$
在 $t<0$ 时递减，在 $t>0$ 时递增。
\textbf{它不是在所有实数上都递增，也不是都递减。}

所以先看 $g(x)$ 的正负：
\[
\begin{array}{c|c|c|c}
 x\text{ 的范围} & t=x-1 & f(t)=t^2 & (x-1)^2\\ \hline
 x<1 & t<0 & \text{随 }t\text{ 增大而减小} & \text{减}\\
 x>1 & t>0 & \text{随 }t\text{ 增大而增大} & \text{增}
\end{array}
\]
\textbf{眼前一亮：不只要确定内层在哪段单调，还要看它的结果会不会跨过外层“掉头”的位置。}
本题的转折点 $t=0$，换回 $x$ 就是 $x=1$。

\medskip
\end{explanationbox}
\clearpage
\begin{explanationbox}
\textbf{\textcolor{CExplanation}{五、另一个意外收获：极大、极小也能直接判断}}
\begin{center}\begin{tikzpicture}[font=\small,>=stealth,line join=round]
\draw[gray!40,rounded corners=3pt] (0,0) rectangle (14.2,11.3);
\node at (7.1,10.700000000000001) {例2：外层递减，内层最低点变成整体最高点};

\begin{scope}[shift={(2,7.3)},x=1.2cm,y=.55cm]
\draw[->,gray] (-1.5,0)--(3.6,0) node[right]{$x$};\draw[->,gray] (0,-1.6)--(0,4.5) node[above]{$t$};
\draw[blue,very thick,domain=-1:3,samples=100] plot(\x,{(\x-1)^2-1});\fill[orange!85!black] (1,-1) circle(2pt);\draw[gray,dotted] (1,-1)--(1,3.7);
\end{scope}
\begin{scope}[shift={(2,2.1)},x=1.2cm,y=1.0cm]
\draw[->,gray] (-1.5,0)--(3.6,0) node[right]{$x$};\draw[->,gray] (0,0)--(0,3.7) node[above]{$H$};
\draw[violet,very thick,domain=-1.3:3.3,samples=150] plot(\x,{pow(3,1-(\x-1)^2)});\fill[orange!85!black] (1,3) circle(2pt);\draw[gray,dotted] (1,0)--(1,3);\node[below=5pt] at (1,0){$1$};
\end{scope}
\node[align=left,anchor=west] at (8.3,8.1){$g(x)=(x-1)^2-1$\\[12pt]$x=1$：内层最小 $-1$\\[12pt]$g'(x)=2(x-1)$};
\node[align=left,anchor=west] at (8.3,3.9){$H(x)=3^{-g(x)}$\\[12pt]外层 $3^{-t}$ 递减\\[12pt]内层的谷，对应整体的峰\\[12pt]$H(1)=3$ 为最大值};

\node[text width=12.8cm,align=center] at (7.1,.65) {两图共享相同的 $x$ 刻度；纵轴分别为内层 $t$ 与整体函数值 $H$。};
\end{tikzpicture}\end{center}



若外层递增，内层的最大值、最小值所在的位置仍分别给出整体的最大值、最小值。
若外层递减，\textbf{内层最小的地方，整体反而最大；内层最大的地方，整体反而最小。}
前提是这些最大、最小值确实能取到，并且都在题目讨论的区间内。

例如 $g(x)=(x-1)^2-1$ 在 $x=1$ 取最小值 $-1$，
而 $f(t)=\left(\frac13\right)^t$ 递减，于是
\[
 y=\left(\frac13\right)^{(x-1)^2-1}
 \quad\text{在 }x=1\text{ 取最大值 }3.
\]
不必展开，更不必求导。

\medskip
\textbf{\textcolor{CExplanation}{六、做题的四步顺序}}
\begin{enumerate}[leftmargin=*, itemsep=0.35em]
\item \textbf{先找能取哪些 $x$：}根号下不能是负数，分母不能为零，对数里面必须大于零。
\item \textbf{再拆里外：}设 $t=g(x)$，把原式写成 $y=f(t)$。
\item \textbf{划分区间：}找内层的转折位置，也找使 $g(x)$ 跨过外层转折位置的 $x$。
\item \textbf{逐段判断：}确认两层在所讨论范围内各自增减，再用同增异减；答案写成 $x$ 的区间。
\end{enumerate}

\textbf{考场快记：}先定义域，后划范围；先里后外；每遇一次“减”就反向一次。
\end{explanationbox}
